% Title page and imprint (verso). Version and date come from main.tex.
\begin{titlepage}
\begin{tikzpicture}[remember picture, overlay]
  \node[anchor=north east, inner sep=0] at ([xshift=-22mm,yshift=-24mm]current page.north east) {\fclogo[0.5]};
\end{tikzpicture}
\vspace*{38mm}
\noindent{\color{fcverm}\rule{1.6cm}{2.2pt}}\par\vspace{14pt}
\noindent{\fcblack\fontsize{34}{39}\selectfont\color{fcink}Fuzzy Clustering\\and Data Analysis\\in Python\par}
\vspace{14pt}
\noindent{\fcmedium\fontsize{16}{20}\selectfont\color{fcbrand}The {\usefont{T1}{IBMPlexMono-TLF}{m}{n}fuzzyclust} toolbox\par}
\vspace{40mm}
\noindent{\fcsemi\fontsize{15}{18}\selectfont\color{fcink}J\'anos Abonyi\par}
\vspace{4pt}
\noindent{\fclightface\fontsize{11}{14}\selectfont\color{fcgray}University of Pannonia, Veszpr\'em\par}
\vfill
\noindent{\fcmedium\small\color{fcgray}Version \fcversion\enspace\textperiodcentered\enspace\fcdate\par}
\end{titlepage}
\newpage
\thispagestyle{empty}
{\small\raggedright
\vspace*{\fill}
\noindent\textbf{Fuzzy Clustering and Data Analysis in Python. The \texttt{fuzzyclust} toolbox}\\
J\'anos Abonyi, University of Pannonia, Veszpr\'em, Hungary\\
Version \fcversion, \fcdate

\bigskip
\noindent\textcopyright~2026 J\'anos Abonyi. This book is licensed under the Creative Commons Attribution 4.0 International License (CC BY 4.0, \url{https://creativecommons.org/licenses/by/4.0/}): it may be copied, redistributed, and adapted for any purpose, provided the source is credited. The \texttt{fuzzyclust} software described in the book is released under the MIT License.

\bigskip
\noindent Software: \url{https://github.com/ComplexSystemsIntelligence/fuzzyclust}\\
Documentation: \url{https://complexsystemsintelligence.github.io/fuzzyclust}\\
Installation: \texttt{pip install fuzzyclust}

\bigskip
\noindent The research summarized in this book was supported by the National Research, Development and Innovation Office of Hungary (NKFIH) through the OTKA 143482 project.

\bigskip
\noindent Typeset with \LaTeX\ in Libertinus Serif, Raleway, and IBM Plex Mono. Every figure that shows a computed result was produced by an example script of the package. The cover shows a Gustafson--Kessel partition computed with \texttt{fuzzyclust}: each sample is drawn in the blend of its three membership colors, and the lines are contours of the memberships.
\par}
